% Source PDF pages 71--74; manual pages 2-42--2-45.
\subsection{Longitude, Latitude, and Altitude Rates}\label{sec:position-rates}

TAOS computes the longitude, latitude, and altitude rates in the function
\texttt{compute\_state\_rates}. The corresponding output variables are
\begin{longtable}{@{}>{\ttfamily}p{0.17\textwidth}p{0.19\textwidth}p{0.56\textwidth}@{}}
  longdt  & $\dot\lambda$ & Longitude rate.\\
  latgcdt & $\dot\delta_{gc}$ & Geocentric latitude rate.\\
  latgddt & $\dot\delta_{gd}$ & Geodetic latitude rate.\\
  altdt   & $\dot h$ & Altitude rate, or rate of climb.\\
\end{longtable}

\taossubhead{Longitude Rate}

Longitude $\lambda$ is defined by
\begin{equation}
  \sin\lambda
  =\frac{\vec r\mathbin{\cdot}\hat y_{\oplus}}
  {\sqrt{
    \left(\vec r\mathbin{\cdot}\hat x_{\oplus}\right)^2
    +\left(\vec r\mathbin{\cdot}\hat y_{\oplus}\right)^2
  }},
  \label{eq:longitude-sine}
\end{equation}
\begin{equation}
  \cos\lambda
  =\frac{\vec r\mathbin{\cdot}\hat x_{\oplus}}
  {\sqrt{
    \left(\vec r\mathbin{\cdot}\hat x_{\oplus}\right)^2
    +\left(\vec r\mathbin{\cdot}\hat y_{\oplus}\right)^2
  }},
  \label{eq:longitude-cosine}
\end{equation}
so that
\begin{equation}
  \lambda
  =\tan^{-1}
    \left[
      \frac{\vec r\mathbin{\cdot}\hat y_{\oplus}}
           {\vec r\mathbin{\cdot}\hat x_{\oplus}}
    \right].
  \label{eq:longitude-arctangent}
\end{equation}
This definition is common to the geocentric and geodetic coordinate systems.
Differentiating it gives
\begin{equation}
  \dot\lambda
  =\frac{
    \left(\vec V_{\oplus}\mathbin{\cdot}\hat y_{\oplus}\right)
    \left(\vec r\mathbin{\cdot}\hat x_{\oplus}\right)
    -
    \left(\vec V_{\oplus}\mathbin{\cdot}\hat x_{\oplus}\right)
    \left(\vec r\mathbin{\cdot}\hat y_{\oplus}\right)
  }{
    \left(\vec r\mathbin{\cdot}\hat x_{\oplus}\right)^2
    +\left(\vec r\mathbin{\cdot}\hat y_{\oplus}\right)^2
  },
  \label{eq:longitude-rate-component-form}
\end{equation}
or, equivalently,
\begin{equation}
  \dot\lambda
  =\frac{
    \vec V_{\oplus}\mathbin{\cdot}
    \left(-\sin\lambda\,\hat x_{\oplus}
          +\cos\lambda\,\hat y_{\oplus}\right)
  }{
    \sqrt{
      \left(\vec r\mathbin{\cdot}\hat x_{\oplus}\right)^2
      +\left(\vec r\mathbin{\cdot}\hat y_{\oplus}\right)^2
    }
  }.
  \label{eq:longitude-rate-unit-vector-form}
\end{equation}
The denominator is the projection of the position vector onto the equatorial plane.
If it becomes zero, longitude rate is undefined and TAOS sets it to zero.

\taossubhead{Geocentric Latitude Rate}

From the local geocentric horizon geometry in
\cref{fig:local-geocentric-horizon}, the geocentric latitude rate is
\begin{equation}
  \dot\delta_{gc}
  =\frac{\vec V_{\oplus}\mathbin{\cdot}\hat x_{gc}}
         {\lVert\vec r\rVert}.
  \label{eq:geocentric-latitude-rate}
\end{equation}

\taossubhead{Geodetic Latitude Rate}

The geodetic latitude rate cannot be isolated as directly. Multiplying
\cref{eq:geodetic-equatorial-projection} by
$(1-e^2)\tan\delta_{gd}$ and subtracting
\cref{eq:geodetic-position-z} gives
\begin{equation}
  \tilde x(1-e^2)\tan\delta_{gd}
  -\vec r\mathbin{\cdot}\hat z_{\oplus}
  =-h e^2\sin\delta_{gd}.
  \label{eq:geodetic-latitude-rate-start}
\end{equation}
Differentiation with respect to time yields
\begin{equation}
  \begin{aligned}
    &\dot{\tilde x}(1-e^2)\tan\delta_{gd}
    +\dot\delta_{gd}\,\tilde x(1-e^2)\sec^2\delta_{gd}
    -\vec V_{\oplus}\mathbin{\cdot}\hat z_{\oplus}\\
    &\qquad=-\dot h\,e^2\sin\delta_{gd}
    -\dot\delta_{gd}\,h e^2\cos\delta_{gd}.
  \end{aligned}
  \label{eq:geodetic-latitude-rate-differentiated}
\end{equation}
Solving for $\dot\delta_{gd}$ gives
\begin{equation}
  \dot\delta_{gd}
  =\frac{
    \left\{
      \dot z\cos\delta_{gd}
      -\left[
        \dot{\tilde x}(1-e^2)
        +\dot h\,e^2\cos\delta_{gd}
       \right]\sin\delta_{gd}
    \right\}\cos\delta_{gd}
  }{
    \tilde x(1-e^2)+h e^2\cos^3\delta_{gd}
  },
  \label{eq:geodetic-latitude-rate-general}
\end{equation}
where
\begin{equation}
  \tilde x=(N+h)\cos\delta_{gd},
  \label{eq:equatorial-projection-for-rate}
\end{equation}
\begin{equation}
  \dot{\tilde x}
  =\vec V_{\oplus}\mathbin{\cdot}
   \left(
     \cos\lambda\,\hat x_{\oplus}
     +\sin\lambda\,\hat y_{\oplus}
   \right),
  \label{eq:equatorial-projection-rate}
\end{equation}
\begin{equation}
  \dot z=\vec V_{\oplus}\mathbin{\cdot}\hat z_{\oplus},
  \label{eq:ecfc-z-position-rate}
\end{equation}
and
\begin{equation}
  \dot h=-\vec V_{\oplus}\mathbin{\cdot}\hat z_{gd}.
  \label{eq:altitude-rate}
\end{equation}
Substituting these relationships into
\cref{eq:geodetic-latitude-rate-general} and simplifying gives
\begin{equation}
  \dot\delta_{gd}
  =\frac{\vec V_{\oplus}\mathbin{\cdot}\hat x_{gd}}
  {N\left(
       \dfrac{1-e^2}{1-e^2\sin^2\delta_{gd}}
     \right)+h}.
  \label{eq:geodetic-latitude-rate-normal-radius}
\end{equation}
Finally, substituting the ellipsoid normal radius $N$ from
\cref{eq:ellipsoid-normal-distance} gives
\begin{equation}
  \dot\delta_{gd}
  =\frac{\vec V_{\oplus}\mathbin{\cdot}\hat x_{gd}}
  {h+
    \dfrac{R_{\oplus}(1-e^2)}
    {\left(1-e^2\sin^2\delta_{gd}\right)^{3/2}}}.
  \label{eq:geodetic-latitude-rate-final}
\end{equation}
The equation is singular when its denominator is zero. In that condition TAOS sets
the geodetic latitude rate equal to the geocentric value.

\taossubhead{Altitude Rate and Acceleration}

The altitude rate is given by \cref{eq:altitude-rate}. Differentiating it gives the
vertical acceleration
\begin{equation}
  \ddot h
  =-\vec a_{\oplus}\mathbin{\cdot}\hat z_{gd}
   -\vec V_{\oplus}\mathbin{\cdot}\dot{\hat z}_{gd},
  \label{eq:altitude-acceleration}
\end{equation}
where differentiation of the geodetic down unit vector gives
\begin{equation}
  \dot{\hat z}_{gd}
  =\frac{d}{dt}
  \left[
    -\cos\delta_{gd}
      \left(
        \cos\lambda\,\hat x_{\oplus}
        +\sin\lambda\,\hat y_{\oplus}
      \right)
    -\sin\delta_{gd}\,\hat z_{\oplus}
  \right],
  \label{eq:geodetic-down-unit-vector-derivative-definition}
\end{equation}
which expands to
\begin{equation}
  \begin{aligned}
  \dot{\hat z}_{gd}={}&
    \left[
      \sin\delta_{gd}
      \left(
        \cos\lambda\,\hat x_{\oplus}
        +\sin\lambda\,\hat y_{\oplus}
      \right)
      -\cos\delta_{gd}\,\hat z_{\oplus}
    \right]\dot\delta_{gd}\\
  &+\cos\delta_{gd}
    \left(
      \sin\lambda\,\hat x_{\oplus}
      -\cos\lambda\,\hat y_{\oplus}
    \right)\dot\lambda.
  \end{aligned}
  \label{eq:geodetic-down-unit-vector-derivative}
\end{equation}
The vertical acceleration is required by one of the guidance rules discussed in
\cref{sec:control-variable-solution}.

\taossubhead{Dynamic Pressure and Mach Number Rates}

The second derivative of dynamic pressure and the first derivative of Mach number
are required by several guidance rules. Dynamic pressure is defined as
\begin{equation}
  q=0.5\,\rho\lVert\vec V_{w\oplus}\rVert^2.
  \label{eq:dynamic-pressure-rate-definition}
\end{equation}
Under the derivative approximations used by TAOS, differentiating twice gives
\begin{equation}
  \ddot q
  =0.5\,\ddot\rho\lVert\vec V_{w\oplus}\rVert^2
   +2\dot\rho\lVert\vec V_{w\oplus}\rVert
      \lVert\dot{\vec V}_{w\oplus}\rVert
   +\rho\lVert\dot{\vec V}_{w\oplus}\rVert^2,
  \label{eq:dynamic-pressure-second-derivative}
\end{equation}
where
\begin{equation}
  \dot\rho
  =\frac{\partial h}{\partial t}\frac{\partial\rho}{\partial h}
  =\dot h\frac{\partial\rho}{\partial h},
  \label{eq:atmospheric-density-rate}
\end{equation}
and
\begin{equation}
  \ddot\rho
  =\frac{\partial^2 h}{\partial t^2}\frac{\partial\rho}{\partial h}
  =\ddot h\frac{\partial\rho}{\partial h}.
  \label{eq:atmospheric-density-second-rate}
\end{equation}
Terms involving $\lVert\ddot{\vec V}_{w\oplus}\rVert$ and
$\partial^2\rho/\partial h^2$ are assumed small,
$\lVert\dot{\vec V}_{w\oplus}\rVert$ is approximated by
$\lVert\vec a_{\oplus}\rVert$, and the atmospheric density derivative is obtained
numerically from the selected atmosphere model.

Mach number is
\begin{equation}
  M=\frac{\lVert\vec V_{w\oplus}\rVert}{c},
  \label{eq:mach-number-rate-definition}
\end{equation}
and its time derivative is
\begin{equation}
  \dot M
  =\frac{
    \lVert\dot{\vec V}_{w\oplus}\rVert c
    -\lVert\vec V_{w\oplus}\rVert\dot c
  }{c^2},
  \label{eq:mach-number-rate}
\end{equation}
where
\begin{equation}
  \dot c
  =\frac{\partial h}{\partial t}\frac{\partial c}{\partial h}
  =\dot h\frac{\partial c}{\partial h}.
  \label{eq:speed-of-sound-rate}
\end{equation}
Again, $\lVert\dot{\vec V}_{w\oplus}\rVert$ is approximated by
$\lVert\vec a_{\oplus}\rVert$, and the altitude derivative of the speed of sound is
computed numerically from the atmosphere model.
